\chapter{1994年旧科目组卷（文）}

\section{单选题}

\begin{problem}
点 \((0,5)\) 到直线 \(y=2x\) 的距离是（\quad）

\choices
    {\(\frac{5}{2}\)}
    {\(\sqrt{5}\)}
    {\(\frac{3}{2}\)}
    {\(\frac{\sqrt{5}}{2}\)}
\end{problem}

\begin{answer}
B
\end{answer}

\begin{solution}
直线 \(y=2x\) 可化为 \(2x-y=0\). 因为点到直线的距离为 \(\frac{|2\cdot 0-5|}{\sqrt{2^{2}+(-1)^{2}}}=\frac{5}{\sqrt{5}}=\sqrt{5}\)，所以选 B.
\end{solution}

\begin{problem}
如果方程 \(x^{2}+ky^{2}=2\) 表示焦点在 \(y\) 轴上的椭圆，那么实数 \(k\) 的取值范围是（\quad）

\choices
    {\((0,+\infty)\)}
    {\((0,2)\)}
    {\((1,+\infty)\)}
    {\((0,1)\)}
\end{problem}

\begin{answer}
D
\end{answer}

\begin{solution}
方程表示椭圆时必有 \(k>0\)，此时方程化为 \(\frac{x^{2}}{2}+\frac{y^{2}}{\frac{2}{k}}=1\). 因为焦点在 \(y\) 轴上，所以 \(y\) 轴方向的分母较大，即 \(\frac{2}{k}>2\)，解得 \(0<k<1\)，故选 D.
\end{solution}

\begin{problem}
\(\displaystyle\lim_{n\to\infty}\frac{\mathrm C_n^0+\mathrm C_n^1+\mathrm C_n^2+\cdots+\mathrm C_n^n}{1+2+2^{2}+\cdots+2^{n}}=\)（\quad）

\choices
    {\(0\)}
    {\(\frac{1}{2}\)}
    {\(1\)}
    {\(2\)}
\end{problem}

\begin{answer}
B
\end{answer}

\begin{solution}
分子是二项式系数之和，即 \(\mathrm C_n^0+\mathrm C_n^1+\cdots+\mathrm C_n^n=2^{n}\)；分母是等比数列之和，即 \(1+2+\cdots+2^{n}=2^{n+1}-1\). 因为 \(\frac{2^{n}}{2^{n+1}-1}=\frac{1}{2-2^{-n}}\)，所以极限为 \(\frac{1}{2}\)，故选 B.
\end{solution}

\begin{problem}
设 \(\theta\) 是第二象限的角，则必有（\quad）

\choices
    {\(\tan\frac{\theta}{2}>\cot\frac{\theta}{2}\)}
    {\(\tan\frac{\theta}{2}<\cot\frac{\theta}{2}\)}
    {\(\sin\frac{\theta}{2}>\cos\frac{\theta}{2}\)}
    {\(\sin\frac{\theta}{2}<\cos\frac{\theta}{2}\)}
\end{problem}

\begin{answer}
A
\end{answer}

\begin{solution}
设 \(\theta=2k\pi+\alpha\)，其中 \(k\in\Z\)，\(\frac{\pi}{2}<\alpha<\pi\)，则 \(\frac{\theta}{2}=k\pi+\frac{\alpha}{2}\)，\(\frac{\pi}{4}<\frac{\alpha}{2}<\frac{\pi}{2}\). 因为正切以 \(\pi\) 为周期，所以 \(\tan\frac{\theta}{2}=\tan\frac{\alpha}{2}>1\)，\(\cot\frac{\theta}{2}=\frac{1}{\tan(\frac{\theta}{2})}<1\)，必有 \(\tan\frac{\theta}{2}>\cot\frac{\theta}{2}\). 而 \(\frac{\theta}{2}\) 可在第一或第三象限，\(\sin\frac{\theta}{2}\) 与 \(\cos\frac{\theta}{2}\) 的大小关系不固定，故选 A.
\end{solution}

\begin{problem}
若直线 \(x+ay+2=0\) 和 \(2x+3y+1=0\) 互相垂直，则 \(a=\)（\quad）

\choices
    {\(-\frac{2}{3}\)}
    {\(-\frac{3}{2}\)}
    {\(\frac{2}{3}\)}
    {\(\frac{3}{2}\)}
\end{problem}

\begin{answer}
A
\end{answer}

\begin{solution}
两直线分别以 \(\left( 1,a \right)\) 和 \(\left( 2,3 \right)\) 为法向量. 因为两直线互相垂直，所以法向量也互相垂直，即 \(1\cdot 2+a\cdot 3=0\)，解得 \(a=-\frac{2}{3}\)，故选 A.
\end{solution}

\begin{problem}
某种细菌在培养过程中，每\(20\text{ 分钟}\)分裂一次（\(1\text{ 个}\)分裂为\(2\text{ 个}\)）. 经过\(3\text{ 小时}\)，这种细菌由\(1\text{ 个}\)可繁殖成（\quad）

\choices
    {\(511\text{ 个}\)}
    {\(512\text{ 个}\)}
    {\(1023\text{ 个}\)}
    {\(1024\text{ 个}\)}
\end{problem}

\begin{answer}
B
\end{answer}

\begin{solution}
\(3\text{ 小时}\)即\(180\text{ 分钟}\)，共分裂 \(180\div 20=9\) 次. 因为每次分裂后数量变为原来的 \(2\) 倍，所以从\(1\text{ 个}\)出发，经过 \(9\) 次分裂后细菌数为 \(2^{9}=512\)，故选 B.
\end{solution}

\begin{problem}
在下列函数中，以 \(\frac{\pi}{2}\) 为周期的函数是（\quad）

\choices
    {\(y=\sin 2x+\cos 4x\)}
    {\(y=\sin 2x\cos 4x\)}
    {\(y=\sin 2x+\cos 2x\)}
    {\(y=\sin 2x\cos 2x\)}
\end{problem}

\begin{answer}
D
\end{answer}

\begin{solution}
\(y=\sin 2x\cos 2x=\frac{1}{2}\sin 4x\)，而 \(\sin 4x\) 的最小正周期是 \(\frac{\pi}{2}\)，所以该函数以 \(\frac{\pi}{2}\) 为周期. 前三项中都含有周期为 \(\pi\) 的 \(\sin 2x\) 或 \(\cos 2x\)，其和或乘积的周期不是 \(\frac{\pi}{2}\)，故选 D.
\end{solution}

\begin{problem}
已知正六棱台的上、下底面边长分别为 \(2\) 和 \(4\)，高为 \(2\)，则其体积为（\quad）

\choices
    {\(32\sqrt{3}\)}
    {\(28\sqrt{3}\)}
    {\(24\sqrt{3}\)}
    {\(20\sqrt{3}\)}
\end{problem}

\begin{answer}
B
\end{answer}

\begin{solution}
边长为 \(s\) 的正六边形面积为 \(6\cdot\frac{\sqrt{3}}{4}s^{2}=\frac{3\sqrt{3}}{2}s^{2}\)，所以上、下底面积分别为 \(6\sqrt{3}\) 和 \(24\sqrt{3}\). 因为棱台体积为 \(\frac{h}{3}\left( S_{1}+S_{2}+\sqrt{S_{1}S_{2}} \right)\)，代入 \(h=2\) 得 \(\frac{2}{3}\left( 6\sqrt{3}+24\sqrt{3}+12\sqrt{3} \right)=28\sqrt{3}\)，故选 B.
\end{solution}

\begin{problem}
使 \(\left( \sqrt{6}-\sqrt{2}\i \right)^{n}\) 是纯虚数的最小自然数 \(n=\)（\quad）

\choices
    {\(3\)}
    {\(4\)}
    {\(5\)}
    {\(6\)}
\end{problem}

\begin{answer}
A
\end{answer}

\begin{solution}
\(\sqrt{6}-\sqrt{2}\i=2\sqrt{2}\left( \frac{\sqrt{3}}{2}-\frac{1}{2}\i \right)=2\sqrt{2}\left( \cos\left( -\frac{\pi}{6} \right)+\i\sin\left( -\frac{\pi}{6} \right) \right)\). 因为其 \(n\) 次方的辐角为 \(-\frac{n\pi}{6}\)，为纯虚数时实部为零，即 \(\cos\left( -\frac{n\pi}{6} \right)=0\)，所以 \(\frac{n\pi}{6}=\frac{\pi}{2}+k\pi\)，\(k\in\Z\)，即 \(n=3+6k\). 满足条件的最小自然数是 \(3\)，故选 A.
\end{solution}

\begin{problem}
有甲、乙、丙三项任务，甲需 \(2\text{ 人}\) 承担，乙、丙各需 \(1\text{ 人}\) 承担，从 \(10\text{ 人}\) 中选派 \(4\text{ 人}\) 承担这三项任务，不同的选法共有（\quad）

\choices
    {\(1260\text{ 种}\)}
    {\(2025\text{ 种}\)}
    {\(2520\text{ 种}\)}
    {\(5040\text{ 种}\)}
\end{problem}

\begin{answer}
C
\end{answer}

\begin{solution}
先从 \(10\text{ 人}\) 中选出承担甲任务的 \(2\text{ 人}\)，有 \(\mathrm C_{10}^{2}\) 种；再从余下 \(8\text{ 人}\) 中选 \(1\text{ 人}\)承担乙任务，有 \(\mathrm C_{8}^{1}\) 种；最后从余下 \(7\text{ 人}\) 中选 \(1\text{ 人}\)承担丙任务，有 \(\mathrm C_{7}^{1}\) 种. 因为 \(\mathrm C_{10}^{2}\mathrm C_{8}^{1}\mathrm C_{7}^{1}=45\cdot 8\cdot 7=2520\)，所以共有 \(2520\text{ 种}\)，故选 C.
\end{solution}

\begin{problem}
对于直线 \(m\)，\(n\) 和平面 \(\alpha\)，\(\beta\)，\(\alpha\perp\beta\) 的一个充分条件是（\quad）

\choices
    {\(m\perp n\)，\(m\parallel\alpha\)，\(n\parallel\beta\)}
    {\(m\perp n\)，\(\alpha\cap\beta=m\)，\(n\subset\alpha\)}
    {\(m\parallel n\)，\(n\perp\beta\)，\(m\subset\alpha\)}
    {\(m\parallel n\)，\(m\perp\alpha\)，\(n\perp\beta\)}
\end{problem}

\begin{answer}
C
\end{answer}

\begin{solution}
因为 \(m\parallel n\)，\(n\perp\beta\)，所以 \(m\perp\beta\)，又 \(m\subset\alpha\)，所以 \(\alpha\perp\beta\)，故选项 C 的条件充分. 反例检验其余选项不充分：选项 A 中 \(\alpha\) 与 \(\beta\) 可以平行；选项 B 只给出 \(n\subset\alpha\)，\(n\perp m\)，得不到 \(n\perp\beta\)；选项 D 中由 \(m\parallel n\)，\(m\perp\alpha\)，\(n\perp\beta\)反而推出 \(\alpha\parallel\beta\)，与 \(\alpha\perp\beta\) 矛盾. 所以选 C.
\end{solution}

\begin{problem}
设函数 \(f(x)=1-\sqrt{1-x^{2}}\)（\(-1\leqslant x\leqslant 0\)），则函数 \(y=f^{-1}(x)\) 的图象是（\quad）

\choices
{\choicebitmap[width=0.15\paperwidth]{img/1994/old_subject_liberal/q12_optA.png}}
{\choicebitmap[width=0.15\paperwidth]{img/1994/old_subject_liberal/q12_optB.png}}
{\choicebitmap[width=0.15\paperwidth]{img/1994/old_subject_liberal/q12_optC.png}}
{\choicebitmap[width=0.15\paperwidth]{img/1994/old_subject_liberal/q12_optD.png}}
\end{problem}

\begin{answer}
B
\end{answer}

\begin{solution}
由 \(y=1-\sqrt{1-x^{2}}\) 得 \(\sqrt{1-x^{2}}=1-y\)，其中 \(0\leqslant y\leqslant 1\)，平方得 \(x^{2}=2y-y^{2}\)，又 \(-1\leqslant x\leqslant 0\)，故反函数为 \(y=-\sqrt{2x-x^{2}}\)（\(0\leqslant x\leqslant 1\)），其图象过点 \((0,0)\)，\((1,-1)\) 且 \(y\leqslant 0\). 端点回代可排除 A，C. 又反函数图象位于圆 \((x-1)^2+y^2=1\) 上，在原点处的切线竖直，可排除原点处切线水平的 D. 所以选 B.
\end{solution}

\begin{problem}
已知过球面上 \(A\)，\(B\)，\(C\) 三点的截面和球心的距离等于球半径的一半，且 \(AB=BC=CA=2\)，则球面面积是（\quad）

\choices
    {\(\frac{16}{9}\pi\)}
    {\(\frac{8}{3}\pi\)}
    {\(4\pi\)}
    {\(\frac{64}{9}\pi\)}
\end{problem}

\begin{answer}
D
\end{answer}

\begin{solution}
设球半径为 \(R\)，则截面到球心的距离为 \(\frac{R}{2}\). 边长为 \(2\) 的正三角形的截面圆半径为 \(r=\frac{2}{\sqrt{3}}\)，由 \(R^{2}=\left( \frac{R}{2} \right)^{2}+r^{2}\) 得 \(\frac{3}{4}R^{2}=\frac{4}{3}\)，解得 \(R=\frac{4}{3}\). 回代检验满足 \(R^{2}-\left( \frac{R}{2} \right)^{2}=r^{2}\). 故球面面积 \(S=4\pi R^{2}=\frac{64}{9}\pi\). 所以选 D.
\end{solution}

\begin{problem}
如果函数 \(y=\sin 2x+a\cos 2x\) 的图象关于直线 \(x=-\frac{\pi}{8}\) 对称，那么 \(a=\)（\quad）

\choices
    {\(\sqrt{2}\)}
    {\(-\sqrt{2}\)}
    {\(1\)}
    {\(-1\)}
\end{problem}

\begin{answer}
D
\end{answer}

\begin{solution}
对称轴处取最值，故 \(y=\sin 2x+a\cos 2x\) 在 \(x=-\frac{\pi}{8}\) 处取最值. 由导数检验亦同：\(y'=2\cos 2x-2a\sin 2x\)，在 \(x=-\frac{\pi}{8}\) 处 \(y'=0\) 得 \(\cos\left( -\frac{\pi}{4} \right)-a\sin\left( -\frac{\pi}{4} \right)=0\)，即 \(1+a=0\)，解得 \(a=-1\). 代入得 \(y=\sin 2x-\cos 2x=\sqrt{2}\sin\left( 2x-\frac{\pi}{4} \right)\)，在 \(x=-\frac{\pi}{8}\) 处确取最小值 \(-\sqrt{2}\)，与辅助角算法一致. 所以选 D.
\end{solution}

\begin{problem}
定义在 \((-\infty,+\infty)\) 上的任意函数 \(f(x)\) 都可以表示成一个奇函数 \(g(x)\) 和一个偶函数 \(h(x)\) 之和. 如果 \(f(x)=\lg(10^{x}+1)\)（\(x\in(-\infty,+\infty)\)），那么（\quad）

\choices
    {\(g(x)=x\)，\(h(x)=\lg(10^{x}+10^{-x}+2)\)}
    {\(g(x)=\frac{1}{2}[\lg(10^{x}+1)+x]\)，\(h(x)=\frac{1}{2}[\lg(10^{x}+1)-x]\)}
    {\(g(x)=\frac{x}{2}\)，\(h(x)=\lg(10^{x}+1)-\frac{x}{2}\)}
    {\(g(x)=-\frac{x}{2}\)，\(h(x)=\lg(10^{x}+1)+\frac{x}{2}\)}
\end{problem}

\begin{answer}
C
\end{answer}

\begin{solution}
设 \(f(x)=g(x)+h(x)\)，其中 \(g(x)\) 为奇函数，\(h(x)\) 为偶函数，则 \(f(-x)=-g(x)+h(x)\)，解得 \(g(x)=\frac{1}{2}[f(x)-f(-x)]\)，\(h(x)=\frac{1}{2}[f(x)+f(-x)]\). 由 \(f(-x)=\lg(10^{-x}+1)=\lg(10^{x}+1)-x\) 得 \(g(x)=\frac{x}{2}\)，\(h(x)=\lg(10^{x}+1)-\frac{x}{2}\). 奇偶性逐项检验排除 A，B，D. 所以选 C.
\end{solution}

\section{填空题}

\begin{problem}
抛物线 \(y^{2}=8-4x\) 的准线方程是\fillinblank{}.
\end{problem}

\begin{answer}
\(x=3\)
\end{answer}

\begin{solution}
由 \(y^{2}=-4\left( x-2 \right)\) 知顶点为 \((2,0)\)，\(2p=4\)，\(p=2\)，开口向左，故准线为 \(x=2+\frac{p}{2}=3\). 化成标准式回代顶点检验一致. 所以准线方程是 \(x=3\).
\end{solution}

\begin{problem}
在 \(\left( x+m \right)^{7}\)（\(m\in\N\)）的展开式中，\(x^{6}\) 的系数是 \(x^{5}\) 的系数与 \(x^{4}\) 的系数的等差中项，则 \(m=\)\fillinblank{}.
\end{problem}

\begin{answer}
\(1\)
\end{answer}

\begin{solution}
通项中 \(x^{6}\)，\(x^{5}\)，\(x^{4}\) 的系数依次为 \(\mathrm C_{7}^{1}m=7m\)，\(\mathrm C_{7}^{2}m^{2}=21m^{2}\)，\(\mathrm C_{7}^{3}m^{3}=35m^{3}\). 由等差中项得 \(2\times 21m^{2}=7m+35m^{3}\)，\(m\in\N\) 为正整数，故 \(m\ne 0\)，两边除以 \(7m\) 得 \(6m=1+5m^{2}\)，解得 \(m=1\) 或 \(m=\frac{1}{5}\)，舍去非自然数解. 代入 \(m=1\) 回代检验 \(7\)，\(21\)，\(35\) 成等差数列. 所以 \(m=1\).
\end{solution}

\begin{problem}
若 \(\frac{5}{4}\pi<\theta<\frac{3}{2}\pi\)，\(\sin 2\theta=a\)，则 \(\cos\theta-\sin\theta\) 的值是\fillinblank{}.
\end{problem}

\begin{answer}
\(\sqrt{1-a}\)
\end{answer}

\begin{solution}
设 \(t=\cos\theta-\sin\theta\)，则 \(t^{2}=1-\sin 2\theta=1-a\). 当 \(\frac{5}{4}\pi<\theta<\frac{3}{2}\pi\) 时 \(\cos\theta<0\)，\(\sin\theta<0\)，取 \(\theta=\frac{4}{3}\pi\) 反例检验得 \(t=\frac{-1+\sqrt{3}}{2}>0\)，故 \(t\) 取正根. 所以所求值为 \(\sqrt{1-a}\).
\end{solution}

\begin{problem}
设圆锥底面圆周上两点 \(A\)，\(B\) 间的距离为 \(2\)，圆锥顶点到直线 \(AB\) 的距离为 \(\sqrt{3}\)，\(AB\) 和圆锥的轴的距离为 \(1\)，则该圆锥的体积为\fillinblank{}.
\end{problem}

\begin{answer}
\(\frac{2\sqrt{2}}{3}\pi\)
\end{answer}

\begin{solution}
设底面中心为 \(O\)，\(M\) 为 \(AB\) 的中点，则 \(OM=1\)，\(AM=1\)，故底面半径 \(r=\sqrt{OM^{2}+AM^{2}}=\sqrt{2}\). 设顶点为 \(S\)，则 \(SM=\sqrt{3}\)，且 \(SO\perp OM\)，\(SO\perp AB\)，所以高 \(h=SO=\sqrt{SM^{2}-OM^{2}}=\sqrt{2}\). 故体积 \(V=\frac{1}{3}\pi r^{2}h=\frac{2\sqrt{2}}{3}\pi\).
\end{solution}

\begin{problem}
在测量某物理量的过程中，因仪器和观察的误差，使得 \(n\) 次测量分别得到 \(a_{1}\)，\(a_{2}\)，\(\cdots\)，\(a_{n}\)，共 \(n\) 个数据. 我们规定所测量物理量的“最佳近似值” \(a\) 是这样一个量：与其他近似值比较，\(a\) 与各数据的差的平方和最小. 依此规定，从 \(a_{1}\)，\(a_{2}\)，\(\cdots\)，\(a_{n}\) 推出的 \(a=\)\fillinblank{}.
\end{problem}

\begin{answer}
\(\frac{1}{n}\left( a_{1}+a_{2}+\cdots+a_{n} \right)\)
\end{answer}

\begin{solution}
记平方和为 \(S=\left( a-a_{1} \right)^{2}+\left( a-a_{2} \right)^{2}+\cdots+\left( a-a_{n} \right)^{2}\)，则 \(S=na^{2}-2a\left( a_{1}+a_{2}+\cdots+a_{n} \right)+\left( a_{1}^{2}+a_{2}^{2}+\cdots+a_{n}^{2} \right)=n\left( a-\frac{1}{n}\left( a_{1}+a_{2}+\cdots+a_{n} \right) \right)^{2}\) 再加与 \(a\) 无关的常数. 求导算法亦得 \(2na-2\left( a_{1}+a_{2}+\cdots+a_{n} \right)=0\). 故当 \(a=\frac{1}{n}\left( a_{1}+a_{2}+\cdots+a_{n} \right)\) 时 \(S\) 最小.
\end{solution}

\section{解答题}

\begin{problem}
求函数 \(y=\frac{\sin3x\sin^{3}x+\cos3x\cos^{3}x}{\cos^{2}2x}+\sin2x\) 的最小值.
\end{problem}

\begin{solution}
因为 \(\cos2x\neq 0\)，由积化和差公式 \(\sin3x\sin x=\frac{\cos2x-\cos4x}{2}\)，\(\cos3x\cos x=\frac{\cos2x+\cos4x}{2}\)，所以
\[ \begin{gathered}
\sin3x\sin^{3}x+\cos3x\cos^{3}x=\frac{1}{2}\left[\left(\cos2x-\cos4x\right)\sin^{2}x+\left(\cos2x+\cos4x\right)\cos^{2}x\right] \\
=\frac{1}{2}\left(\cos2x+\cos2x\cos4x\right)=\cos^{3}2x
\end{gathered} \]
所以 \(y=\cos2x+\sin2x\)，即 \(y=\sqrt{2}\sin\left(2x+\frac{\pi}{4}\right)\). 当 \(\sin\left(2x+\frac{\pi}{4}\right)=-1\)，即 \(x=-\frac{3\pi}{8}+k\pi\)（\(k\in\Z\)）时，\(\cos2x=-\frac{\sqrt{2}}{2}\neq 0\)，符合定义域，所以 \(y\) 取最小值 \(-\sqrt{2}\).
\end{solution}

\begin{problem}
已知函数 \(f(x)=\lg x\)（\(x\in\R^{+}\)）. 若 \(x_{1}\)，\(x_{2}\in\R^{+}\)，判断 \(\frac{1}{2}\left[f\left(x_{1}\right)+f\left(x_{2}\right)\right]\) 与 \(f\left(\frac{x_{1}+x_{2}}{2}\right)\) 的大小，并加以证明.
\end{problem}

\begin{solution}
因为 \(f\left(x_{1}\right)+f\left(x_{2}\right)=\lg x_{1}+\lg x_{2}=\lg\left(x_{1}x_{2}\right)\)，由均值不等式 \(x_{1}x_{2}\leq\left(\frac{x_{1}+x_{2}}{2}\right)^{2}\)，当且仅当 \(x_{1}=x_{2}\) 时取等号. 因为常用对数底大于 \(1\)，所以 \(\lg\left(x_{1}x_{2}\right)\leq 2\lg\left(\frac{x_{1}+x_{2}}{2}\right)\)，即 \(\frac{1}{2}\left[f\left(x_{1}\right)+f\left(x_{2}\right)\right]\leq f\left(\frac{x_{1}+x_{2}}{2}\right)\)，当且仅当 \(x_{1}=x_{2}\) 时取等号.
\end{solution}

\begin{problem}
如图，已知 \(A_{1}B_{1}C_{1}-ABC\) 是正三棱柱，\(D\) 是 \(AC\) 中点.
\begin{enumerate}
\item 证明 \(AB_{1}\parallel\)平面\(DBC_{1}\)；
\item 假设 \(AB_{1}\perp BC_{1}\)，\(BC=2\)，求线段 \(AB_{1}\) 在侧面 \(B_{1}BCC_{1}\) 上的射影长.
\end{enumerate}
\bitmapfigure[width=0.5\linewidth]{img/1994/old_subject_liberal/q23_fig1.png}
\end{problem}

\begin{solution}
\begin{enumerate}
\item 连结 \(B_{1}C\) 交 \(BC_{1}\) 于 \(E\)，连结 \(DE\). 因为 \(A_{1}B_{1}C_{1}-ABC\) 是正三棱柱，所以四边形 \(B_{1}BCC_{1}\) 是矩形，对角线互相平分，故 \(B_{1}E=EC\). 在 \(\triangle AB_{1}C\) 中，因为 \(AD=DC\)，所以 \(DE\parallel AB_{1}\). 又 \(AB_{1}\) 不在平面 \(DBC_{1}\) 内，\(DE\) 在平面 \(DBC_{1}\) 内，所以 \(AB_{1}\parallel\)平面\(DBC_{1}\).
\item 作 \(AF\perp BC\)，垂足为 \(F\). 因为 \(A_{1}B_{1}C_{1}-ABC\) 是正三棱柱，侧棱垂直于底面，所以平面 \(ABC\perp\)平面\(B_{1}BCC_{1}\)；又 \(AF\perp BC\)，\(BC\) 为交线，所以 \(AF\perp\)平面\(B_{1}BCC_{1}\). 连结 \(B_{1}F\)，则 \(B_{1}F\) 是 \(AB_{1}\) 在平面 \(B_{1}BCC_{1}\) 内的射影. 因为 \(AF\perp\)平面\(B_{1}BCC_{1}\)，所以 \(AF\perp BC_{1}\)；又 \(AB_{1}\perp BC_{1}\)，所以 \(BC_{1}\perp\)平面\(AB_{1}F\)，故 \(BC_{1}\perp B_{1}F\). 在矩形 \(B_{1}BCC_{1}\) 中，\(\angle B_{1}BF=\angle BCC_{1}=90^{\circ}\)；由 \(B_{1}F\perp BC_{1}\) 得 \(\angle FB_{1}B\) 与 \(\angle BFB_{1}\) 互余，且 \(\angle C_{1}BC\) 与 \(\angle BFB_{1}\) 互余，故 \(\angle FB_{1}B=\angle C_{1}BC\)，所以 \(\triangle B_{1}BF\sim\triangle BCC_{1}\)，得 \(\frac{B_{1}B}{BC}=\frac{BF}{CC_{1}}\). 因为 \(F\) 为正三角形 \(ABC\) 的 \(BC\) 边中点，\(BC=2\)，所以 \(BF=1\)；又 \(CC_{1}=B_{1}B\)，故 \(B_{1}B^{2}=BF\cdot BC=2\). 于是 \(B_{1}F^{2}=B_{1}B^{2}+BF^{2}=3\)，即 \(B_{1}F=\sqrt{3}\). 所以线段 \(AB_{1}\) 在侧面 \(B_{1}BCC_{1}\) 上的射影长为 \(\sqrt{3}\).
\end{enumerate}
\end{solution}

\begin{problem}
已知直角坐标平面上一点 \(Q(2,0)\) 和圆 \(C\)，\(x^{2}+y^{2}=1\)，动点 \(M\) 到圆 \(C\) 的切线长与 \(|MQ|\) 的比等于 \(\sqrt{2}\). 求动点 \(M\) 的轨迹方程，说明它表示什么曲线.
\bitmapfigure[width=0.5\linewidth]{img/1994/old_subject_liberal/q24_fig1.png}
\end{problem}

\begin{solution}
设切点为 \(N\)，点 \(M\) 的坐标为 \((x,y)\). 因为圆 \(C\) 半径 \(|ON|=1\)，所以 \(|MN|^{2}=|MO|^{2}-|ON|^{2}=x^{2}+y^{2}-1\). 由 \(|MN|=\sqrt{2}|MQ|\) 得 \(x^{2}+y^{2}-1=2\left[\left(x-2\right)^{2}+y^{2}\right]\). 整理得 \(x^{2}+y^{2}-8x+9=0\)，即 \(\left(x-4\right)^{2}+y^{2}=7\). 经检验，在圆 \(\left(x-4\right)^{2}+y^{2}=7\) 上 \(x\geqslant 4-\sqrt{7}\)，\(x^{2}+y^{2}-1=8x-10\geqslant 22-8\sqrt{7}>0\)，切线长有意义，坐标适合方程的点都满足 \(|MN|=\sqrt{2}|MQ|\)，故该方程为所求轨迹方程，它表示圆心为 \((4,0)\)、半径为 \(\sqrt{7}\) 的圆.
\end{solution}

\begin{problem}
设数列 \(\{a_{n}\}\) 的前 \(n\) 项和为 \(S_{n}\)，若对所有自然数 \(n\)，都有 \(S_{n}=\frac{n\left(a_{1}+a_{n}\right)}{2}\)，证明 \(\{a_{n}\}\) 是等差数列.
\end{problem}

\begin{solution}
当 \(n\geqslant 2\) 时，由题设 \(S_{n}=\frac{n\left(a_{1}+a_{n}\right)}{2}\)，\(S_{n-1}=\frac{\left(n-1\right)\left(a_{1}+a_{n-1}\right)}{2}\)，所以 \(a_{n}=S_{n}-S_{n-1}=\frac{n\left(a_{1}+a_{n}\right)}{2}-\frac{\left(n-1\right)\left(a_{1}+a_{n-1}\right)}{2}\). 将上式中 \(n\) 换成 \(n+1\)，得 \(a_{n+1}=\frac{\left(n+1\right)\left(a_{1}+a_{n+1}\right)}{2}-\frac{n\left(a_{1}+a_{n}\right)}{2}\). 于是 \(a_{n+1}-a_{n}=\left(S_{n+1}-S_{n}\right)-\left(S_{n}-S_{n-1}\right)=S_{n+1}-2S_{n}+S_{n-1}\)，即 \(a_{n+1}-a_{n}=\frac{\left(n+1\right)\left(a_{1}+a_{n+1}\right)}{2}-n\left(a_{1}+a_{n}\right)+\frac{\left(n-1\right)\left(a_{1}+a_{n-1}\right)}{2}\). 上式两边同乘以 \(2\)，\(a_{1}\) 的系数为 \(\left(n+1\right)-2n+\left(n-1\right)=0\)，故 \(\left(n-1\right)a_{n+1}-2\left(n-1\right)a_{n}+\left(n-1\right)a_{n-1}=0\). 因为 \(n\geqslant 2\)，所以 \(a_{n+1}-a_{n}=a_{n}-a_{n-1}\). 故 \(a_{n+1}-a_{n}=a_{n}-a_{n-1}=\cdots=a_{2}-a_{1}\)，\(\{a_{n}\}\) 是等差数列.
\end{solution}
